\documentclass[10pt,a4paper]{amsart}
\usepackage[margin=19mm]{geometry}
\usepackage{amsmath,amssymb,mathtools}
\usepackage[scr=rsfs,cal=euler]{mathalfa}
\usepackage{xfrac}
\usepackage[unicode,hidelinks,bookmarks=false]{hyperref}
\newcommand{\ie}{i.e.\ }
\newcommand{\cf}{cf.\ }
\newcommand{\resp}{respectively\ }
\newcommand{\Subsection}{\subsection}
\providecommand{\nobreakdash}{\nobreak\hbox{-}}
\setlength{\emergencystretch}{2em}
\multlinegap0pt
\usepackage{xcolor}
\usepackage{algorithm}\usepackage{algpseudocode}
\usepackage{amssymb}
\usepackage{mathtools}
\usepackage{tcolorbox}
\usepackage{tikz}
\usepackage{enumerate}
\newtheorem{proposition}{Proposition}
\newcommand{\dual}[1]{\left\langle {#1} \right\rangle}
\providecommand{\paragraph}{}
\title{Adaptive Accelerated Mirror Descent in Primal and Dual Spaces}
\author{Zeyi Xu, Long Chen}
\date{Source: arXiv:2606.02787v1 (CC BY 4.0)}
\begin{document}
\maketitle

\noindent\emph{Source and licence.} Adaptive Accelerated Mirror Descent in Primal and Dual Spaces. Version arXiv:2606.02787v1 (CC BY 4.0), distributed by arXiv under the Creative Commons Attribution 4.0 International licence, http://creativecommons.org/licenses/by/4.0/, as recorded in the arXiv licence metadata for this exact version and shown on the abstract page https://arxiv.org/abs/2606.02787v1. Full licence notice: \texttt{licenses/arxiv-2606.02787-CC-BY-4.0.txt}.\par\smallskip

\noindent\textit{Editorial scope.} Counted unit: the article's proposition proposition (source0.tex lines 1431--1433). The article's complete original proof of it is delivered with it (source0.tex lines 1434--1704, one \texttt{proof} environment of 271 lines). No mathematics is written, repaired or re-derived: the statement and the proof are the article's own lines, in the article's own notation, and no retained line is substituted or re-flowed.  No further source-local context is needed: every cross-reference of every retained line resolves inside the two delivered spans.  One declared adaptation: exactly 2 ancillary strings inside the retained spans are omitted (image-inclusion commands and, where the source repeats a label, the repeated label definitions) and no other line differs from the source; the figure environments, with their placement, captions and labels, are kept, and the package records the omitted strings in fidelity receipts bound to the affected bodies.  No retained line contains a citation, so the wrapper carries no bibliography and no bibliography entry.  The retained lines contain the article's own diagram code, which the wrapper typesets with the standard tikz packages; no image file is delivered and the package declares no asset dependency.  Editorial formatting only, in addition: an AMS \texttt{amsart} preamble that restates the article's own declarations and macros used by the retained lines, namely the environments \texttt{proposition} on separate counters, together with the standard packages those lines need.  The retained environments print in this document as Proposition 1 in order of appearance, whereas the article prints the same items with the numbering of its own full document.  The complete native source of the article as distributed in the arXiv e-print for this exact version is bundled unchanged as \texttt{sources/201925/source0.tex}.
\begin{proposition}
    The function $f$ has the anisotropic descent property of $\frac{1}{L}$ with respect to the preconditioner function $\phi$ if and only if $\phi^*$ is $L$-relatively smooth with respect to $f^*$, i.e., $D_{\phi^*}(\chi, \eta)\leq L D_{f^*}(\chi, \eta)$ for all $\chi, \eta\in V^*$.
\end{proposition}

\begin{proof}
Throughout, let $x,y\in\mathbb{R}^n$ be arbitrary and set
\[
\chi:=\nabla f(x),\qquad \eta:=\nabla f(y),
\qquad x=\nabla f^*(\chi),\qquad y=\nabla f^*(\eta).
\]
Assume $L=1$ (as explained via epi-scaling). We prove the equivalence between

\smallskip
\noindent\emph{(i)} $\phi^*$ is $1$-relatively smooth with respect to $f^*$, i.e.
\begin{equation}\label{eq:rel-smooth-1}
D_{\phi^*}(\eta, \chi)\le D_{f^*}(\eta, \chi),\qquad \forall\,\chi,\eta\in V^*,
\end{equation}
and

\smallskip
\noindent\emph{(ii)} $f$ has the anisotropic descent property of $1$ with respect to $\phi$, i.e.
\begin{equation}\label{eq:aniso-descent-1}
D_f(x,y)\le D_{\phi}(x-y^+,\,y-y^+),
\end{equation}
By definition,
\[
y-y^+ = \nabla \phi^*(\eta), \quad \eta = \nabla f(y), \quad \eta = \nabla \phi(y-y^+).
\]
%since $\nabla\phi$ and $\nabla\phi^*$ are inverse maps on their respective domains.

{\bf Part: $\eqref{eq:aniso-descent-1}\ \Rightarrow\ \eqref{eq:rel-smooth-1}$.}


Fix $x,y\in\mathbb{R}^n$ and set
\[
h(z):=D_f(z,y) = f(z) - f(y) - \dual{\nabla f(y), z - y},\qquad z\in\mathbb{R}^n.
\]
Then $h$ is convex, $h(z)\ge 0$ for all $z$, and $y = \arg\min h(z)$.
Moreover, $h$ differs from $f$ by an affine term, hence they share the same Bregman divergence.

We use the anisotropic descent property to construct an upper approximation $h_{\phi}(z)$ of $h(z)$ at $z = x$. We use $h_{\phi}$ to define a gradient descent step at $x$. Then calculate the minimizer of $h_{\phi}$ and compare the difference to get \eqref{eq:rel-smooth-1}.

\medskip

\paragraph{\em Step 1: Construct an upper approximation of $h$.}
By the Bregman decomposition with base point $x$,
\begin{equation}\label{eq:h-bregman}
h(z)
=
h(x)+\langle \nabla h(x),z-x\rangle + D_h(z,x),
=
h(x)+\langle \nabla h(x),z-x\rangle + D_f(z,x),
\end{equation}
where
\[
\nabla h(x)=\nabla f(x)-\nabla f(y)=\chi - \eta.
\]

By assumption \eqref{eq:aniso-descent-1}, we have the upper bound
\[
D_f(z,x)\ \le\ D_\phi(z-x^+,x-x^+),
\]
where 
$x^+:=x-\nabla\phi^*(\chi), \nabla\phi(x-x^+)=\chi.$
Plugging this into \eqref{eq:h-bregman}, we obtain the upper model
\[
h(z)\ \le\ h_\phi(z)
:= h(x)+\langle \nabla h(x),z-x\rangle + D_\phi(z-x^+,x-x^+).
\]

\begin{figure}[htbp]
\begin{center}

\caption{Upper approximation of $D_f$ at $x$.}
\label{fig:upper}
\end{center}
\end{figure}

\paragraph{\em Step 2: Compute the minimizer $z^\star$ of $h_\phi$.}
We have
\[
\nabla h_\phi(z)
=
(\chi-\eta)+\nabla\phi(z-x^+)-\nabla\phi(x-x^+)
= \nabla\phi(z-x^+) - \eta.
\]
Setting $\nabla h_\phi(z^\star)=0$ gives
\[
\nabla\phi(z^\star-x^+)=\eta,
\qquad\text{so}\qquad
z^\star-x^+ = \nabla\phi^*(\eta).
\]
Using $x^+=x-\nabla\phi^*(\chi)$, we conclude
\begin{equation}\label{eq:zstar}
z^\star
- x=
\nabla\phi^*(\eta) -\nabla\phi^*(\chi).
\end{equation}


\paragraph{\em Step 3: Evaluate $h_\phi(x)-h_\phi(z^\star)$.}
First, $h_\phi(x)=h(x)$.
For $z=z^\star$, using \eqref{eq:zstar},
\[
z^\star-x
=
-\nabla\phi^*(\chi)+\nabla\phi^*(\eta),
\qquad
z^\star-x^+
=\nabla\phi^*(\eta),
\qquad
x-x^+=\nabla\phi^*(\chi),
\]
we obtain
\begin{align*}
h_\phi(z^\star)
&=
h(x)+\langle \chi-\eta,z^\star-x\rangle
+D_\phi(z^\star-x^+,x-x^+)\\
&=
h(x)+\langle \chi-\eta,\nabla\phi^*(\eta)-\nabla\phi^*(\chi)\rangle
+D_\phi\bigl(\nabla\phi^*(\eta),\nabla\phi^*(\chi)\bigr).
\end{align*}
Hence
\begin{equation}\label{eq:delta-def}
\Delta
:=
h_\phi(x)-h_\phi(z^\star)
=
-\langle \chi-\eta,\nabla\phi^*(\eta)-\nabla\phi^*(\chi)\rangle
- D_\phi\bigl(\nabla\phi^*(\eta),\nabla\phi^*(\chi)\bigr).
\end{equation}

%\medskip
%\paragraph{\em Step 5: Simplify $\Delta$ in terms of $D_{\phi^*}$.}
By conjugacy of Bregman divergences,
\[
D_\phi\bigl(\nabla\phi^*(\eta),\nabla\phi^*(\chi)\bigr)
=
D_{\phi^*}(\chi,\eta).
\]
Moreover, the symmetrization identity for $\phi^*$ gives
\[
D_{\phi^*}(\eta,\chi)+D_{\phi^*}(\chi,\eta)
=
\langle \nabla\phi^*(\eta)-\nabla\phi^*(\chi),\,\eta-\chi\rangle.
\]
Therefore,
\begin{align*}
\Delta
&=
-\langle \chi-\eta,\nabla\phi^*(\eta)-\nabla\phi^*(\chi)\rangle
- D_\phi\bigl(\nabla\phi^*(\eta),\nabla\phi^*(\chi)\bigr)\\
&=
-\langle \chi-\eta,\nabla\phi^*(\eta)-\nabla\phi^*(\chi)\rangle
- D_{\phi^*}(\chi,\eta)\\
&=
D_{\phi^*}(\eta,\chi)+D_{\phi^*}(\chi,\eta)-D_{\phi^*}(\chi,\eta)
= D_{\phi^*}(\eta,\chi).
\end{align*}

\medskip
\paragraph{\em Step 4: Relate $\Delta$ to $D_f(x,y)$.}
Since $h\ge h_\phi$ pointwise and $h(y)=\min_z h(z)$, we have
\[
h(x)-h(y)\ \ge\ h(x)-\min_z h_\phi(z)\ \ge\ h(x)-h_\phi(z^\star)=\Delta.
\]
Based on Fig. \ref{fig:upper}, we obtain
\[
D_{f^*}(\eta,\chi) = D_f(x,y)
= h(x)-h(y)
\ge h(x)-h_\phi(z^\star)
= D_{\phi^*}(\eta,\chi),
\]
which is exactly \eqref{eq:rel-smooth-1}.

\medskip
%------------------------------------------------------------
{\bf Part: $\eqref{eq:rel-smooth-1}\ \Rightarrow\ \eqref{eq:aniso-descent-1}$.}

The argument is analogous to Part~1, but carried out in the dual space for the dual objective $f^*$.

Fix $x\in\mathbb{R}^n$ and let $\chi=\nabla f(x)$. Define
\[
h(v):=D_{f^*}(v,\chi)
= f^*(v)-f^*(\chi)-\langle \nabla f^*(\chi),v-\chi\rangle,
\qquad v\in\mathbb{R}^n.
\]
Then $h(v)\ge 0$ for all $v$, with $\chi=\arg\min_v h(v)$. Since $h$ differs from $f^*$ only by an affine function, they share the same Bregman divergence:
\[
D_h(u,v)=D_{f^*}(u,v), \qquad \forall\,u,v.
\]
Moreover,
\[
\nabla h(\eta)=\nabla f^*(\eta)-\nabla f^*(\chi)=y-x.
\]

\paragraph{\em Step 1: Construct a lower approximation of $h$.}
By the Bregman decomposition of $h$ at $\eta$, we have
\[
h(v)
= h(\eta)+\langle \nabla h(\eta),v-\eta\rangle + D_{f^*}(v,\eta).
\]
Using the relative smoothness assumption \eqref{eq:rel-smooth-1},
$D_{f^*}(v,\eta)\ge D_{\phi^*}(v,\eta)$, we obtain the lower model
\begin{equation}\label{eq:key-lower-bound}
h(v)\ \ge\ h_{\phi^*}(v)
:= h(\eta)+\langle \nabla h(\eta),v-\eta\rangle + D_{\phi^*}(v,\eta).
\end{equation}

\begin{figure}[htbp]
\begin{center}

\caption{Lower approximation of $D_{f^*}$ at $\eta$.}
\label{fig:Dflower}
\end{center}
\end{figure}

\paragraph{\em Step 2: Compute the minimizer $v^\star$ of $h_{\phi^*}$.}
The gradient of the nonconstant part of $h_{\phi^*}$ is
\[
\nabla_v\Big(\langle \nabla h(\eta),v-\eta\rangle + D_{\phi^*}(v,\eta)\Big)
= \nabla h(\eta)+\nabla\phi^*(v)-\nabla\phi^*(\eta).
\]
Setting this to zero yields
\begin{equation}\label{eq:vstar}
\nabla\phi^*(v^\star)
= \nabla\phi^*(\eta)-\nabla h(\eta)
= \nabla\phi^*(\eta)-(y-x)=x-y^+,
\end{equation}
that is,
\[
v^\star=\nabla\phi(x-y^+).
\]

\paragraph{\em Step 3: Evaluate the height of the lower model.}
Since $h_{\phi^*}(\eta)=h(\eta)$, we compute
\[
h_{\phi^*}(\eta)-h_{\phi^*}(v^\star)
= -\langle \nabla h(\eta),v^\star-\eta\rangle - D_{\phi^*}(v^\star,\eta).
\]
Let
\[
a:=x-y^+,\qquad b:=y-y^+.
\]
Then $y-x=b-a$, $v^\star=\nabla\phi(a)$, and $\eta=\nabla\phi(b)$. Hence
\[
-\langle \nabla h(\eta),v^\star-\eta\rangle
= \langle b-a,\,\nabla\phi(b)-\nabla\phi(a)\rangle.
\]
Using conjugacy,
$D_{\phi^*}(\nabla\phi(a),\nabla\phi(b))=D_\phi(b,a)$, and the symmetrized
Bregman identity, we obtain
\begin{equation}\label{eq:height-phi}
h_{\phi^*}(\eta)-h_{\phi^*}(v^\star)
= D_\phi(a,b)
= D_\phi(x-y^+,\,y-y^+).
\end{equation}

\paragraph{\em Step 4: Compare with the height of $h$.}
Since $h\ge h_{\phi^*}$ pointwise and $h(\chi)=\min_v h(v)$,
\[
h(\eta)-h(\chi)
\le h_{\phi^*}(\eta)-\min_v h_{\phi^*}(v)
= h_{\phi^*}(\eta)-h_{\phi^*}(v^\star).
\]
Therefore,
\[
D_f(x,y)
= D_{f^*}(\eta,\chi)
\le D_\phi(x-y^+,\,y-y^+),
\]
which is exactly the anisotropic descent property \eqref{eq:aniso-descent-1}.

\end{proof}

\end{document}
